\scanpage{15}
\mspage{\handwritten{14}}

catégorie \inserted{additive(et même quasi-abélienne)}\margin{\handwritten{$=$ karoubienne $+$ additive}} avec $\otimes$ qui a toutes les vertus, sauf d'être abélienne, et sauf d'admettre des \uline{Hom} i\deleted{\unsure{j|n}}nternes. Pour trouver ces derniers, on peut si on veut paraphraser la construction de la catégorie des motifs, \unsure{en} \deleted{introduisant un \unsure{motif}} \inserted{\handwritten{en}} considérant l'objet $\inserted{\handwritten{\Lambda}(-1)}$ de $\handwritten{\widetilde{\underline B}}$ conoyau de $p^{\handwritten{\unsure{*|\bullet}}}:e\longrightarrow S$, où $p:S^1\longrightarrow e$ est l'application du cercle sur un point. C'est un objet de poids $1$ (le poids est défini par le degré \deleted{\unsure{d}} des générateurs, changés de signe). On constate aussitôt que le foncteur $M\mapsto M(-1)=M\otimes\handwritten{\Lambda}(-1)$ de $\handwritten{\widetilde{\underline B}}$ dans lui-même est pleinement \unsure{fidèle}, ce qui permet alors, par une construction formelle standard, de comp\inserted{\handwritten{lé}}ter $\handwritten{\widetilde{\underline B}}$ en une catégorie où $\handwritten{\Lambda}(-1)$ devient inversible. Il se trouve alors que dans cette dernière les $\underline{\mathrm{Hom}}(X,Y)$ internes existent et \deleted{satisfont} sont isomorphes à $\overset{\handwritten{\vee}}X\otimes Y$, ... en fait , on trouve maintenant \handwritten{la} catégorie de\deleted{\unsure{x|s}} \uline{tous} les modules gradués libres de type fini sur $\handwritten{\Lambda}$. Bien sûr, cette catégorie n'est pas encore abélienne pour autant, puisque $\handwritten{\Lambda}$ n'est pas un corps; \deleted{pour}\inserted{si on veut} définir la catégorie abélienne engendrée, on s'attend à trouver la\deleted{s} catégorie desmodules de \deleted{type} \inserted{\handwritten{présentation}} \unsure{finie} sur $\handwritten{\Lambda}$. (\deleted{\unsure{---} est-il noethérien, i.e. le nombre de générateurs est-il fini ? ).}\inserted{\handwritten{qui est abélienne bien que $\Lambda$ ne soit pas \unsure{noethérien|noéthérien}, car $\Lambda$ est $\unsure{\underline{\lim}|\underset{\longrightarrow}{\lim}}$ d'anneaux \unsure{noéth.|noeth.} à hom de transition \uline{plats})}}.
\margin{\handwritten{\unsure{?|??}}}

Cette construction ressemblerait à celle d'une théorie des motifs où on n'aurait pas passé à l'équivalence homologique pour les cycles, mais où on aurait gardé l'équivalence rationelle, - qui est bien trop fine \deleted{\unsure{pou|pour}}à beaucoup d'égards. Pour trouver un équivalent plus étroit, il \unsure{convinnt|convint|convient|convinmt} de reprendre la théorie précédente,on pourrait formuler le même Pb universel que plus haut, mais en \unsure{prenant} l'équivalence homologique (via le graphe) des morphismes de variétés au lieu de l'équivalence homotopique. \deleted{Mais alors il semble que l'on trouve comme catégorie solution celle dont les objets sont les varié}

% archive sheet number: 14
